\documentclass[a4paper,12pt]{article}
\usepackage[utf8]{inputenc}
\usepackage{geometry}
\usepackage{listings}
\usepackage{xcolor}
\usepackage{titlesec}
\usepackage{hyperref}
\geometry{margin=1in}

\titleformat{\section}{\normalfont\Large\bfseries}{\thesection}{1em}{}
\titleformat{\subsection}{\normalfont\large\bfseries}{\thesubsection}{1em}{}

\title{Library Management System Project Report}
\author{Mahla Entezari}
\date{Mar 2023}

\begin{document}

\maketitle

\tableofcontents
\newpage

\section{Introduction}
This project is a simple command-line based Library Management System written in Java. It supports operations for both users and librarians, including borrowing and returning books, managing book inventories, and user registration.

\section{System Overview}
The system is designed with an object-oriented approach and consists of several classes, each representing a component of the library system.

\section{Class Descriptions}

\subsection{Book}
The \texttt{Book} class contains basic information about a book including its name, author, publication year, and ISBN. It provides getter and setter methods to access and modify these attributes.

\subsection{User}
The \texttt{User} class represents a library user. It includes a username, password, and a list of borrowed books. It allows users to rent and return books.

\subsection{Librarian}
The \texttt{Librarian} class holds credentials of library administrators who are authorized to manage the book collection.

\subsection{Library}
This is the core class of the system. It maintains lists of books, users, and librarians and provides functionality to:

\begin{itemize}
    \item Add, remove, update, and search for books
    \item Add, remove, update, and search for users
    \item Add, remove, and search for librarians
    \item Track the quantity of each book
\end{itemize}

\subsection{Main}
The \texttt{Main} class runs the application. It includes the menu-driven interface where users and librarians can interact with the system.

\section{Functionalities}

\subsection{User Operations}
\begin{itemize}
    \item Login and logout
    \item Register as a new user
    \item Borrow books
    \item Return books
\end{itemize}

\subsection{Librarian Operations}
\begin{itemize}
    \item Login and logout
    \item Add new books
    \item Remove books
    \item Increase/decrease quantity of books
\end{itemize}

\section{Sample Code Snippet}
Here is an example of how the system allows a user to borrow a book:
\begin{lstlisting}[language=Java, basicstyle=\ttfamily\small, keywordstyle=\color{blue}]
if (!library.doesBookExist(name, author, year, isbn)) {
    System.err.println("This book doesn't exist!");
} else {
    library.removeBook(name, author, year, isbn);
}
\end{lstlisting}

\section{Conclusion}
The project demonstrates fundamental principles of object-oriented programming. Future enhancements may include persistent storage, a GUI, or user authentication with encryption.

\end{document}
